\documentclass{book}
%\documentclass{book}
%If you're going to print this, perhaps get ride of oneside so that \parts{} will have their own page!
\title{
	\begin{huge}
		\bf{Automorphism Representability of $\R$}
	\end{huge}
	}
\date{December 20, 2025}
\author{Nathanael Chwojko-Srawley}

%\date{} % If you want to specify a date

%  ╔══════════════════════════════════════════════════════════╗
%  ║                         packages                         ║
%  ╚══════════════════════════════════════════════════════════╝

\usepackage[a4paper, total={6in, 8in}]{geometry}
\usepackage[answerdelayed]{exercise}
\usepackage[font=itshape]{quoting}
\usepackage[most]{tcolorbox}         % Makes nice boxes for thms, cor, prop, or anything else.
\usepackage{amsmath, amssymb, amsthm} % Used to write better mathematical expressions (ex. \mathbb{R})
\usepackage{array}                   % \begin{table}{>{\em}c c} -> 1st column italic (bfseries for bold)
\usepackage{bbm}                     % for mathbbm{1}
\usepackage{blindtext}
\usepackage{commutative-diagrams}
\usepackage{dsfont}
\usepackage{dynkin-diagrams}         % for Dynkin diagrams
\usepackage{enumitem}
\usepackage{etoolbox}
\usepackage{euscript}
\usepackage{extarrows}               % Used to compile any UTF-8 Character (and supports Unicode)
\usepackage{fancyhdr}
\usepackage{float}                   % package with ``H'' for figures and tables, ex. Images, tables, etc
\usepackage{graphicx}                % Let's you use images in LaTeX; ex the \includegraphics command.
\usepackage{hhline}
\usepackage{hyperref}
\usepackage{cleveref}
\usepackage{inputenc}                % Used to compile any UTF-8 Character (and supports Unicode)
\usepackage{listings}                % Create nice colour code blocks (customized in preamble)
\usepackage{longtable}               % create tables that extend past one page
\usepackage{makeidx}
\usepackage{mathrsfs}
\usepackage{mathtools}
% \usepackage{microtype}               % makes nice text. Fixes some under-/over- filled box problems.
\usepackage{multicol}                % For multiple columns
\usepackage{pgfplots}
\usepackage{scalerel}
\usepackage{setspace}                % More flexibility on spacing in document.
\usepackage{stackengine}
\usepackage{stackrel}                % for left-right arrows, staking symbol above and bellow
\usepackage{thmtools}                % Customize theorem environment (in preamble).
\usepackage{tikz-cd}                 % for commutative diagrams https://tex.stackexchange.com/questions/419221/how-to-do-the-pushout-with-universal-property
\usepackage{tikz}                    % for commutative diagrams https://tex.stackexchange.com/questions/419221/how-to-do-the-pushout-with-universal-property
\usepackage{titlesec}
\usepackage{wasysym}                 % emoji's!
\usepackage{wrapfig}                 % To wrap text around figure
\usepackage{xcolor}                  % Colour text.

% \usepackage{comment}

% \usepackage{CJKutf8} % for some chinese

% \usepackage{dutchcal} % for lower case mathcal
% \pdfsuppresswarningpagegroup=1

% ╔═════════════════════════════════════════════════════╗
% ║           for figures via inkspace                  ║
% ╚═════════════════════════════════════════════════════╝

% to import figures via: https://github.com/gillescastel/inkscape-figures
\usepackage{import}
\usepackage{pdfpages}
\usepackage{transparent}


% This command is the ONLY command imported via the packages snippet, think of it as a tiny package written out
\newcommand{\incfig}[2][1]{%
    \def\svgwidth{#1\columnwidth}
    \import{./figures/}{#2.pdf_tex}
}

% ╔═══════════════════════════════════════════════════════════╗
% ║                 bibliography and citing                   ║
% ╚═══════════════════════════════════════════════════════════╝

\usepackage{csquotes}
% \usepackage[backend=biber,citestyle=alphabetic,bibstyle=authortitle]{biblatex}
\usepackage[backend=biber]{biblatex}

\addbibresource{bibliography.bib}

% ╔═══════════════════════════════════════════════════════════╗
% ║                    colours (colors)                       ║
% ╚═══════════════════════════════════════════════════════════╝

\definecolor{dkgreen}{rgb}{0,0.6,0}
\definecolor{gray}{rgb}{0.5,0.5,0.5}
\definecolor{mauve}{rgb}{0.58,0,0.82}
\definecolor{lightyellow}{rgb}{1,1,0.6}
\definecolor{reddish}{HTML}{A01A3A}

%  ╔══════════════════════════════════════════════════════════╗
%  ║               title/heading/footer format                ║
%  ╚══════════════════════════════════════════════════════════╝

\pagestyle{fancy}
\setlength\headheight{20pt}
\renewcommand{\sectionmark}[1]{\markright{\thesection.\ #1}}
\renewcommand{\chaptermark}[1]{\markboth{#1}{}}
\fancyfoot[C,CO]{\textcolor{gray}{Nathanael Chwojko-Srawley}}
\fancyfoot[R,RO]{\thepage}{}

\titleformat
{\chapter} % command
[display] % shape
{\bfseries\Huge\itshape} % format
{\thechapter} % label
{0.5ex} % sep
{
    \rule{\textwidth}{1pt}
    \vspace{1ex}
    \centering
}
[
\vspace{-0.5ex}
\rule{\textwidth}{0.3pt}
]

\setlength{\parindent}{0pt} % don't like indentation infront of paragraphs
\setlength{\parskip}{1ex}   %so that there is still space between paragarphs


%  ╔══════════════════════════════════════════════════════════╗
%  ║                       Shortcuts                          ║
%  ╚══════════════════════════════════════════════════════════╝

%\ExerciseLevelInToc

% I like original mathcal better, but need lowercase.
\DeclareFontFamily{OT1}{pzc}{}
\DeclareFontShape{OT1}{pzc}{m}{it}{<-> s * [1.10] pzcmi7t}{}
\DeclareMathAlphabet{\mathpzc}{OT1}{pzc}{m}{it}


% mathbb
\newcommand{\A}{{\mathbb{A}}}
\newcommand{\C}{{\mathbb{C}}}
\newcommand{\F}{{\mathbb{F}}}
\newcommand{\N}{{\mathbb{N}}}
\newcommand{\Q}{{\mathbb{Q}}}
\newcommand{\R}{{\mathbb{R}}}
\newcommand{\V}{{\mathbb{V}}}
\newcommand{\Z}{{\mathbb{Z}}}
\newcommand{\BI}{{\mathbb{I}}}
\renewcommand{\H}{{\mathbb{H}}}
\renewcommand{\P}{{\mathbb{P}}}


% mathcal
% \newcommand{\co}{{\mathfrak{o}}} %mathcal{o} looks like wreath product, so I changed it to mathfrak
\newcommand{\co}{{\mathpzc{o}}}
\newcommand{\CA}{{\mathcal{A}}}
\newcommand{\CB}{{\mathcal{B}}}
\newcommand{\CC}{{\mathcal{C}}}
\newcommand{\CD}{{\mathcal{D}}}
\newcommand{\CF}{{\mathcal{F}}}
\newcommand{\CG}{{\mathcal{G}}}
\newcommand{\CH}{{\mathcal{H}}}
\newcommand{\CI}{{\mathcal{I}}}
\newcommand{\CK}{{\mathcal{K}}}
\newcommand{\CL}{{\mathcal{L}}}
\newcommand{\CM}{{\mathcal{M}}}
\newcommand{\CN}{{\mathcal{N}}}
\newcommand{\CO}{{\mathcal{O}}}
\newcommand{\CQ}{{\mathcal{Q}}}
\newcommand{\CR}{{\mathcal{R}}}
\newcommand{\CS}{{\mathcal{S}}}
\newcommand{\CT}{{\mathcal{T}}}
\newcommand{\CV}{{\mathcal{V}}}
\newcommand{\CZ}{{\mathcal{Z}}}
% EXCPETION:
\newcommand{\CP}{{\mathcal{P}}}
% \newcommand{\CP}{\mathbb{C}\mathrm{P}}
\newcommand{\RP}{\mathbb{R}\mathrm{P}}


%Euscript
\newcommand{\EB}{\EuScript{B}}
\newcommand{\EC}{{\EuScript{C}}}
\newcommand{\EF}{{\EuScript{F}}}
\newcommand{\EL}{{\EuScript{L}}}
\newcommand{\EO}{{\EuScript{O}}}


%mathfrak (lowercase)
\newcommand{\fa}{{\mathfrak{a}}}
\newcommand{\fb}{{\mathfrak{b}}}
\newcommand{\fc}{{\mathfrak{c}}}
\newcommand{\fd}{{\mathfrak{d}}}
\newcommand{\fm}{{\mathfrak{m}}}
\newcommand{\fn}{{\mathfrak{n}}}
\newcommand{\fo}{{\mathfrak{o}}}
\newcommand{\fp}{{\mathfrak{p}}}
\newcommand{\fq}{{\mathfrak{q}}}


%mathfrak (upper case)
\newcommand{\FA}{{\mathfrak{A}}}
\newcommand{\FB}{{\mathfrak{B}}}
\newcommand{\FC}{{\mathfrak{C}}}
\newcommand{\FD}{{\mathfrak{D}}}
\newcommand{\FK}{{\mathfrak{K}}}
\newcommand{\FM}{{\mathfrak{M}}}
\newcommand{\FN}{{\mathfrak{N}}}
\newcommand{\FO}{{\mathfrak{O}}}
\newcommand{\FP}{{\mathfrak{P}}}


%mathscr
\newcommand{\ff}{{\mathscr{F}}}
\newcommand{\SA}{\mathscr{A}}
\newcommand{\SC}{\mathscr{C}}
\newcommand{\SF}{\mathscr{F}}
\newcommand{\SG}{\mathscr{G}}
\newcommand{\SH}{\mathscr{H}}
\newcommand{\SI}{\mathscr{I}}
\newcommand{\SK}{\mathscr{K}}
\newcommand{\SN}{\mathscr{N}}
\newcommand{\SQ}{\mathscr{Q}}
\newcommand{\SR}{\mathscr{R}}
\newcommand{\SV}{\mathscr{V}}
\newcommand{\SZ}{\mathscr{Z}}


% operatorname -- 2 letters (lowercase and upper case)
\newcommand{\ab}{{\operatorname{ab}}}
\newcommand{\ad}{{\operatorname{ad}}}
\newcommand{\ch}{{\operatorname{ch}}}
\newcommand{\ev}{{\operatorname{ev}}}
\newcommand{\id}{{\operatorname{id}}}
\newcommand{\im}{{\operatorname{im}}}
\newcommand{\nm}{{\operatorname{Nm}}}
\newcommand{\op}{{\operatorname{op}}}
\newcommand{\rk}{{\operatorname{rk}}}
\newcommand{\sh}{{\operatorname{sh}}}
\newcommand{\tr}{{\operatorname{tr}}}

\newcommand{\Ab}{{\operatorname{Ab}}}
\newcommand{\Br}{{\operatorname{Br}}}
\newcommand{\Ch}{{\operatorname{Ch}}}
\newcommand{\Cl}{{\operatorname{Cl}}}
\newcommand{\GL}{{\operatorname{GL}}}
\newcommand{\Nm}{{\operatorname{Nm}}}
\newcommand{\SL}{{\operatorname{SL}}}

\renewcommand{\Re}{{\operatorname{Re}}}
\renewcommand{\Im}{{\operatorname{Im}}}


%categories
\newcommand{\Grp}{{\textbf{Grp}}}
\newcommand{\Fld}{{\textbf{Fld}}}
\newcommand{\Sh}{{\textbf{Sh}}}
\newcommand{\Set}{{\textbf{Set}}}
\newcommand{\Mod}{{\textbf{Mod}}}
\newcommand{\PSh}{{\textbf{PSh}}}
\newcommand{\Sch}{{\textbf{Sch}}}
\newcommand{\Top}{{\textbf{Top}}}
\newcommand{\RgSp}{{\textbf{RgSp}}}
\newcommand{\Ring}{{\textbf{Ring}}}


%operatorname -- 3+ letters (lowercase and upper case)
\newcommand{\rad}{{\operatorname{rad}}}
\newcommand{\sep}{{\operatorname{sep}}}
\newcommand{\res}{{\operatorname{res}}}
\newcommand{\sgn}{{\operatorname{sgn}}}
\newcommand{\pre}{{\operatorname{pre}}}
\newcommand{\ord}{{\operatorname{ord}}}
\newcommand{\vol}{{\operatorname{vol}}}
\newcommand{\lcm}{{\operatorname{lcm}}}

\newcommand{\conj}{{\operatorname{conj}}}
\newcommand{\disc}{{\operatorname{disc}}}
\newcommand{\stab}{{\operatorname{stab}}}
\newcommand{\kdim}{{\operatorname{kdim}}}
\newcommand{\diag}{{\operatorname{diag}}}

\newcommand{\trdeg}{\operatorname{trdeg}}
\newcommand{\coker}{{\operatorname{coker}}}
\newcommand{\codim}{{\operatorname{codim}}}

\newcommand{\Ann}{{\operatorname{Ann}}}
\newcommand{\Aut}{{\operatorname{Aut}}}
\newcommand{\Der}{{\operatorname{Der}}}
\newcommand{\Div}{{\operatorname{Div}}}
\newcommand{\Emb}{{\operatorname{Emb}}}
\newcommand{\End}{{\operatorname{End}}}
\newcommand{\Ext}{{\operatorname{Ext}}}
\newcommand{\Gal}{{\operatorname{Gal}}}
\newcommand{\Hom}{{\operatorname{Hom}}}
\newcommand{\Ind}{{\operatorname{Ind}}}
\newcommand{\Inn}{{\operatorname{Inn}}}
\newcommand{\Int}{{\operatorname{int}}}
\newcommand{\Irr}{{\operatorname{Irr}}}
\newcommand{\Jac}{{\operatorname{Jac}}}
\newcommand{\Mor}[1]{\operatorname{Mor}(\textbf{#1})}
\newcommand{\Nil}{{\operatorname{Nil}}}
\newcommand{\Obj}{{\operatorname{Obj}}}
\newcommand{\Out}{{\operatorname{Out}}}
\newcommand{\Pic}{{\operatorname{Pic}\,}}
\newcommand{\Rep}{{\operatorname{Rep}}}
\newcommand{\Res}{{\operatorname{Res}}}
\newcommand{\Spl}{{\operatorname{Spl}}}
\newcommand{\Syl}{{\operatorname{Syl}}}
\newcommand{\Sym}{{\operatorname{Sym}}}
\newcommand{\Tor}{{\operatorname{Tor}}}

\newcommand{\Char}{{\operatorname{char}}}
\newcommand{\Frac}{{\operatorname{Frac}}}
\newcommand{\Frob}{{\operatorname{Frob}}}
\newcommand{\Proj}{{\operatorname{Proj}\,}}
\newcommand{\RMod}{{}_R\operatorname{Mod}}
\newcommand{\SExt}{{\mathcal{E}\emph{xt}}}
\newcommand{\SHom}{{\emph{Hom}}}
\newcommand{\Span}{{\operatorname{span}}}
\newcommand{\Spec}{{\operatorname{Spec}\,}}
\newcommand{\Supp}{{\operatorname{Supp}}}
\newcommand{\Tens}[1]{#1\otimes --}
\newcommand{\fHom}[1]{\operatorname{Hom}(#1, --)}

\newcommand{\mSpec}{{\operatorname{mSpec}}}


% connectors
% \renewcommand{\mid}{\big|}
\newcommand{\sse}{\subseteq}
\newcommand{\subg}{\leqslant}
\newcommand{\supg}{\geqslant}
\newcommand{\ssne}{\subsetneq}
\newcommand{\isosubg}{\lesssim}
\newcommand{\nsse}{\not\subseteq}
\newcommand{\nsubg}{\trianglelefteq}
\newcommand{\nsupg}{\trianglerighteq}
\newcommand{\csubg}{\blacktriangleleft}
\renewcommand{\mid}{\big|}


% Arrows
\newcommand{\rw}{\rightarrow}
\newcommand{\Rw}{\Rightarrow}
\newcommand{\lw}{\leftarrow}
\newcommand{\Lw}{\Leftarrow}
\newcommand{\lrw}{\leftrightarrow}
\newcommand{\LRw}{\Leftrightarrow}
\newcommand{\acton}{\curvearrowright}
\newcommand{\actson}{\curvearrowright}
\newcommand\mapsfrom{\mathrel{\reflectbox{\ensuremath{\mapsto}}}}

% arrows for commutative diagram: create four new angled double-struck arrows
\newcommand\myrot[1]{\mathrel{\rotatebox[origin=c]{#1}{$\Rightarrow$}}}
\newcommand\NEarrow{\myrot{45}}
\newcommand\NWarrow{\myrot{135}}
\newcommand\SWarrow{\myrot{-135}}
\newcommand\SEarrow{\myrot{-45}}


%shorthands
\newcommand{\oln}[1]{{\overline{#1}}}
\renewcommand{\bf}[1]{\textbf{#1}}
\newcommand{\qn}{\quad\newline}


% limits
\newcommand{\Lim}{{\varprojlim}}
\newcommand{\coLim}{{\varinjlim}}
\newcommand{\Dlim}{\lim\limits_{\longrightarrow}}
\newcommand{\coDlim}{\lim\limits_{\longleftarrow}}
\newcommand{\cofHom}[1]{\operatorname{Hom}(--, #1)}


%for saturated ideals in the localization section (I don't like these, I think I'll delete)
\newcommand{\LIm}{{}^e }
\newcommand{\LPre}{{}^c }

%  ╔══════════════════════════════════════════════════════════╗
%  ║                    Custom Commands                       ║
%  ╚══════════════════════════════════════════════════════════╝

% swap varphi and phi
\let\temp\phi
\let\phi\varphi
\let\varphi\temp

\newcommand{\placeholder}{\_\_\_\_\_}

\newcommand{\set}[2]{\left\{#1 \ \middle|\ #2\right\}}

%mod should not have the crazy amount of space to its right!
\renewcommand{\mod}{\,\operatorname{mod}\,}

% dcups
\newcommand{\dcup}{\sqcup}
\newcommand{\Dcup}{\coprod}
% \newcommand{\Dcup}{\bigsqcup}

\newcommand{\nbot}{\mathrel{\rlap{$\bot$}\mkern2mu/}}

%a nice large circle
\newcommand{\tikzcircle}[2][red,fill=black]{\tikz[baseline=-0.5ex]\draw[#1,radius=#2] (0,0) circle ;}%

\makeatletter
\newcommand*\bigcdot{\mathpalette\bigcdot@{.5}}
\newcommand*\bigcdot@[2]{\mathbin{\vcenter{\hbox{\scalebox{#2}{$\m@th#1\bullet$}}}}}
\makeatother

%somehow not a default command
\def\rddots{\mathstrut^{.^{.^{.}}}}

%% new delimiter
\DeclarePairedDelimiter{\ceil}{\lceil}{\rceil}


%  ╔══════════════════════════════════════════════════════════╗
%  ║                     environments                         ║
%  ╚══════════════════════════════════════════════════════════╝


%remember the option: breakable

% create tcolor boxes
\newtcbtheorem[number within=section]{thm}{Theorem}%
{
colback=green!5,
colframe=green!35!black,
fonttitle=\bfseries,
label type=theorem
}{th}

\newtcbtheorem[number within=section, use counter from=thm]{lem}{Lemma}%
{
colback=blue!5,
colframe=black!35!black,
fonttitle=\bfseries,
label type=lemma
}{lm}
\newtcbtheorem[number within=section, use counter from=thm]{prop}{Proposition}%
{
colback=white!5,
colframe=white!35!black,
fonttitle=\bfseries,
label type=proposition
}{pr}
\newtcbtheorem[number within=section, use counter from=thm]{cor}{Corollary}%
{colback=blue!5,
colframe=blue!35!black,
fonttitle=\bfseries,
label type=corollary
}{co}
\newtcbtheorem[number within=section, use counter from=thm]{axiom}{Axiom}%
{colback=black!5,
colframe=yellow!35!black,
fonttitle=\bfseries,
label type=axiom
}{ax}
\newtcbtheorem[number within=section, use counter from=thm]{defn}{Definition}%
{colback=red!5,
colframe=red!35!black,
fonttitle=\bfseries,
label type=definition
}{df}


% for <s-cr> to create \item[$\LRw$] instead of \item
\newenvironment{equivEnumerate}
 {\begin{enumerate}
	}{
\end{enumerate}}

%insure the exercises are properly numbered
\counterwithin{Exercise}{section}
\counterwithin{Answer}{section}


%Custom enviroments
 \newenvironment{note}
 {\begin{tcolorbox}[enhanced, sharp corners, colback=white , borderline={0.2pt}{0pt}{black}]
   \begin{quote}\textbf{Note}:
   }{
   \end{quote}\end{tcolorbox}}

 \newenvironment{intuition}
 {\begin{quote}\textbf{Intuition}:
   }{
   \end{quote}}

\newtcbtheorem[number within=section, use counter from=thm]{titledBox}{}{%
  enhanced,
  sharp corners,
  colback=white,
  colframe=black,
  borderline={0.2pt}{0pt}{black},
  left=2mm,
  right=2mm,
  top=2mm,
  bottom=2mm,
  title style={opacity=1},
  attach title to upper,
  before upper={\textbf{\tcbtitletext}\quad},
  coltitle=black,
  fonttitle=\bfseries,
  separator sign={\quad}
}{box}


%better example and proof environments
\def\exampletext{Example} % If English
\colorlet{colexam}{red!55!black} % Global example color


\newtcbtheorem[number within=section, use counter from=thm]{example}{Example}{% \newtcolorbox[use counter=testexample]{testexamplebox}{%
% Example Frame Start
empty,% Empty previously set parameters
title={\exampletext \thetcbcounter #1},% use \thetcbcounter to access the testexample counter text
% Attaching a box requires an overlay
attach boxed title to top left,
   % Ensures proper line breaking in longer titles
   minipage boxed title,
% (boxed title style requires an overlay)
boxed title style={empty,size=minimal,toprule=0pt,top=4pt,left=3mm,overlay={}},
coltitle=colexam,fonttitle=\bfseries,
before=\par\medskip\noindent,parbox=false,boxsep=0pt,left=3mm,right=0mm,top=2pt,breakable,pad at break=0mm,
   before upper=\csname @totalleftmargin\endcsname0pt, % Use instead of parbox=true. This ensures parskip is inherited by box.
% Handles box when it exists on one page only
overlay unbroken={\draw[colexam,line width=.5pt] ([xshift=-0pt]title.north west) -- ([xshift=-0pt]frame.south west); },
% Handles multipage box: first page
overlay first={\draw[colexam,line width=.5pt] ([xshift=-0pt]title.north west) -- ([xshift=-0pt]frame.south west); },
% Handles multipage box: middle page
overlay middle={\draw[colexam,line width=.5pt] ([xshift=-0pt]frame.north west) -- ([xshift=-0pt]frame.south west); },
% Handles multipage box: last page
overlay last={\draw[colexam,line width=.5pt] ([xshift=-0pt]frame.north west) -- ([xshift=-0pt]frame.south west); }%
}{ex}



\newtcolorbox{ProofTcb}[1][]{%
    empty,
    title={\emph{Proof}: #1},
    attach boxed title to top left,
    minipage boxed title,
    boxed title style={empty,size=minimal,toprule=0pt,top=4pt,left=3mm,overlay={}},
    coltitle=black!55!black,
    fonttitle=\bfseries,
    before=\par\medskip\noindent,
    parbox=false,
    boxsep=0pt,
    left=3mm,
    right=0mm,
    top=2pt,
    breakable,
    pad at break=0mm,
    before upper=\csname @totalleftmargin\endcsname0pt,
    width=\dimexpr\textwidth-3mm\relax,
    overlay unbroken={\draw[black!55!black,line width=.5pt] ([xshift=-0pt]title.north west) -- ([xshift=-0pt]frame.south west); },
    overlay first={\draw[black!55!black,line width=.5pt] ([xshift=-0pt]title.north west) -- ([xshift=-0pt]frame.south west); },
    overlay middle={\draw[black!55!black,line width=.5pt] ([xshift=-0pt]frame.north west) -- ([xshift=-0pt]frame.south west); },
    overlay last={\draw[black!55!black,line width=.5pt] ([xshift=-0pt]frame.north west) -- ([xshift=-0pt]frame.south west); },
    #1
}

\newenvironment{Proof}[1][]
  {\begin{ProofTcb}[#1]}
  {\end{ProofTcb}}


%  ╔══════════════════════════════════════════════════════════╗
%  ║                    for Dynkin diagrams                   ║
%  ╚══════════════════════════════════════════════════════════╝

\def\row#1/#2!{#1_{\IfStrEq{#2}{}{n}{#2}} & \dynkin{#1}{#2}\\}
\newcommand{\tble}[1]{
   \renewcommand*\do[1]{\row##1!}
   \[
      \begin{array}{ll}\docsvlist{#1}\end{array}
   \]
}

%  ╔══════════════════════════════════════════════════════════╗
%  ║                         links                            ║
%  ╚══════════════════════════════════════════════════════════╝

\definecolor{LinkColour}{HTML}{A64A3B} % favourite
% \definecolor{LinkColour}{HTML}{8C3D32} % 2nd place
\hypersetup{
  colorlinks,
  citecolor=black,
  filecolor=magenta,
  linkcolor=LinkColour,
  urlcolor=blue,
}



%  ╔══════════════════════════════════════════════════════════╗
%  ║                          misc                            ║
%  ╚══════════════════════════════════════════════════════════╝

\stackMath %to be able to stack text in math mode
\pgfplotsset{compat=1.18}
\makeindex

%import options: all, bibliography, book-formatting, booksSetting, customEnvironments, dynkin, hw-formatting, language-setup, newcommands, packages, preamble, proofs, settings, tcolorbox, test.svg,
\begin{document}
\maketitle
% \tableofcontents

A group $H$ is \emph{[group] automorphism-representable} if there exists a group $G$ such that
\[
  \Aut_\Grp(G) \cong H
\]
From now on we'll just write $\Aut$ instead of $\Aut_\Grp$.

We shall prove the following:
\begin{thm}{Torsion-Free Divisible Abelian Groups not Automorphism Representable}{torFreeAbGrpNotAutRep}
  Let $H$ be a torsion-free divisible abelian group. Then there does not exist a group $G$ such that
  \[
    \Aut(G) \cong H
  \]
\end{thm}

This result is a generalization of the following simpler result:

\begin{lem}{Rational Are Not Automorphism-representable}{rationalAreNotAutomorphism-representable}
  The rational numbers under addition $(\Q, +)$ are not automorphism representable
\end{lem}

\begin{Proof}
  The key to remember is that finitely generated subgroups of $\Q$ are cyclic.
  For the sake of contradiction say they are so that there exists a $G$ such that:
  \[
    \Aut(G) \cong \Q
  \]
  Then as:
  \[
    \Inn(G) \nsubg \Aut(G) \cong \Q
  \]
  either $\Inn(G)$ is a divisible group or cyclic. If it is is cyclic then it must have infinite order as $\Q$
  is torsion free. Then as:
  \[
    G/Z(G) \cong \Z
  \]
  Elementary group theory tells us $G$ is abelian. But if $G$ is abelian there must be a torsion element
  $x \mapsto x^{-1}$ of order $2$, a contradiction. On the other hand, if $G/Z(G)$ is not finitely generated,
  we can ``cheat'' by choosing any $gZ(G), hZ(G) \in G/Z(G)$ and generating the group $H$ which is cyclic, so
  \[
    \langle xZ(G) \rangle = \langle gZ(G), hZ(G) \rangle
  \]
  Then it is easy to see that:
  \[
    gh = x^nz_1x^mz_2 = z^{n+m}z_1z_2 = z^mz_2x^nz_1 = hg
  \]
  and as $g,h$ were arbitrary, $G$ is abelian giving us the same contradiction.
\end{Proof}

Another important fact: let us make sure that if $G$ is an $\F_2$-vector-space then $\Aut(G)$ cannot be
abelian. The point of this is that we want to abuse the $x \mapsto x^{-1}$ argument, but we can't do that if
$x^{-1} = x$ and hence this is the identity:

\begin{lem}{Vector Space, then Non-abelian}{vecSpthenNon-abelian}
  Assume $x \mapsto x^{-1}$ is the identity on the group $G$. Then $G$ has an $\F_2$-action, hence is a vector space and hence
  $\Aut(G)$ cannot be abelian
\end{lem}

\begin{Proof}
  Say that $x \mapsto x^{-1}$ is the identity. Then every element has a natural $\F_2$-action on it. Hence $G$
  is a vector-space which we know by linear algebra $\Aut(G) = \GL(G)$ is not abelian.
\end{Proof}

We shall be clever with such arguments in a moment.

\subsection*{Central automorphism and cohomology}

We shall explicitly aim to prove that if $A$ is a torsion-free divisible abelian group (like $\Q$, $\Q^n$, or
$\R$), then it is not automorphism representable.For the sake of contradiction let's say $\Aut(G) \cong A$. As
$A$ is abelian, we can conclude all automorphisms are \emph{central}:

\begin{defn}{Central Automorphisms}{centralAutos}
    An automorphism $\alpha \in \Aut(G)$ is called \emph{central} if it induces the identity map on the
    central quotient $G/Z(G)$. Equivalently, for all $x \in G$, $x^{-1}\alpha(x) \in Z(G)$. In other words:
    \[
      \alpha(x) \in xZ(G)
    \]
    so $\alpha$ keeps elements in the cosets.
\end{defn}

We claimed that:
\[
  \Aut_c(G) \cong \Aut(G)
\]
To show this, since $\Aut(G)$ is abelian, the subgroup of inner automorphisms, $\Inn(G) \cong G/Z(G)$, must be
a normal subgroup of $\Aut(G)$. Let $\alpha \in \Aut(G)$ and $c_g \in \Inn(G)$ (where $c_g(x) = gxg^{-1}$).
Since $\Aut(G)$ is abelian, $\alpha$ and $c_g$ commute:
\[
    \alpha \circ c_g = c_g \circ \alpha
\]
Applying this to all $x \in G$:
\begin{align*}
    \alpha(gxg^{-1}) &= g\alpha(x)g^{-1} \\
    \alpha(g)\alpha(x)\alpha(g)^{-1} &= g\alpha(x)g^{-1}
\end{align*}
Since $\alpha$ is surjective, as $x$ varies over $G$, $\alpha(x)$ varies over all of $G$.  Letting $y
= \alpha(x)$, re-arranging we see that:
\[
  g^{-1} \alpha(g) \, y = y \, g^{-1} \alpha(g)
\]
Letting $z = g^{-1}\alpha(g)$ we get:
\[
  zy = yz
\]
Hence, as $y$ varies through all elements of $G$,
\[
  z = g^{-1}\alpha(g) \in Z(G) \qquad \text{i.e.} \qquad \alpha(g) \in gZ(G)
\]
which is exactly the condition of a central automorphism. Thus if $\Aut(G)$ is abelian, then
\[
  \Aut(G) = \Aut_c(G)
\]

We shall abuse this with some cohomology. First, from group theory we know that $\Inn(G) \cong G/Z(G)$ so
that:
\[
  0 \to Z(G) \to G \to \Inn(G) \to 0
\]
is a central extension, where $\Inn(G) \subg A$ is abelian. hence, $G$ is a nilpotent group of order $2$,
namely $[G,G] \subg Z(G)$. This will allow for great cohomological restraints. Recall that $H^2(Z(G),
\Inn(G))$ is in bijective correspondences with all extensions, characterized by the choice of 2-cocycles.
Namely, if $Q = G/Z(G)$ and $Z = Z(G)$, then since $G$ is a central extension of $Z$ by $Q$, elements of $G$ can be identified with pairs $(a, z)$ where $a \in Q$ and $z \in Z$, with the multiplication rule:
\[
(a, z_1) \cdot (b, z_2) = (a+b, z_1 + z_2 + f(a, b))
\]
where $f: Q \times Q \to Z$ is a 2-cocycle (factor set); see~(\textcolor{red}{cite Neukirch and Robinson and
myself}). Note that $f(a,b) \ne f(b,a)$ (or else this group would be abelian and we know this implies
$\Aut(G)$ has to have an order $2$ element), and how we are using the fact that the 2nd term is
the center, or else we would have to write:
\[
(a, z_1) \cdot (b, z_2) = (a+b, a\cdot z_1 + z_2 + f(a, b))
\]
for some action of $a$ on $z_1$, but this action is trivial using the center (ibid on sources).

Now comes the crucial part: as $G$ is nilpotent of class $2$, we have the following powerful theorem for
nilpotent groups: By Warfield (\textit{Nilpotent Groups}, Theorem 5.4), the extension class of $G$ in $H^2(Q,
Z)$ maps to a commutator form $\omega \in \operatorname{Hom}(\Lambda^2 Q, Z)$, namely:
\[
f(x,y) \mapsto  f(x,y) - f(y,x) = \omega(u, v) \in \Hom(\Lambda^2 Q, Z)
\]
In his words:
\begin{thm}{Warfield, Thm 5.4}{warfield54}
  If $A$ and $B$ are Abelian groups, and $H^2(B, A)$ is the group of central extensions of $A$ by $B$, then there is an exact sequence
  \[
  0 \to \operatorname{Ext}(B, A) \to H^2(B, A) \to \operatorname{Hom}(\Lambda^2 B, A) \to 0
  \]
  where the map on the right is given by the commutator form. If $A$ is uniquely 2-divisible, the sequence has
  a canonical splitting. Otherwise, the sequence splits, but not necessarily in a unique way.
\end{thm}
The theorem states that $\omega$ is an \emph{Alternating Bilinear Form}. As it is bilinear, it implies:
\[
  \omega(-u, -v) = (-1)(-1)\omega(u, v) = \omega(u, v)
\]

\subsection{Using this result to eliminate some immediate cases}

The main thins we need is that $A = G/Z(G)$ is at least $2$-divisible, since in that case the sequence splits
and we can write:
\[
  f(x,y) = \frac{1}{2} \left(f(x,y) - f(y,x)\right) + \frac{1}{2} \left(f(x,y) + f(y,x)\right)
\]
If this is the case, then we can define a homomorphism (clearly bijective, hence isomorphism):
\[
  \Psi: G \to G \qquad \Psi(a,b) = (-a,b)
\]
indeed:
\begin{align*}
\Psi((a, z_1) \cdot (b, z_2)) &= \Psi(a+b, z_1+z_2+f(a, b)) \\
&= (-(a+b), z_1+z_2+f(a, b))
\end{align*}
\begin{align*}
\Psi(a, z_1) \cdot \Psi(b, z_2) &= (-a, z_1) \cdot (-b, z_2) \\
&= (-a-b, z_1+z_2+f(-a, -b))
\end{align*}
Since $f(a, b) = f(-a, -b)$, the map preserves the group operation. Since $\Psi^2 = \text{id}$, it is bijective and thus an automorphism.
Note the important of bi-linearity to insure it is well-defined. It is now an order $2$ automorphism,
a contradiction .

Similarly, if $B$ or $A$ is free or even projective (so $Z(G)$ or $G/Z(G)$ are free), then $\Ext(B,A) = \{0\}$
and we get that $f(x,y)$ must be an alternating form, and so we get the same argument.

\subsection{Lie-Algebra case}

Next, it may be that $A$ is not $2$-divisible, but it \emph{is} still $t$-divisible for some $t \in \N$, we employ
the following equivalence of categories with rational nilpotent lie algebras.

\textcolor{red}{here}.

\subsection{Where the ideas stop working}

In the mixed case, say $A = \Q\times \Z$, then I am stuck.



\subsection{Further ideas}


% Now, each $\alpha \in \Aut_c(G)$ gives rise to a derivation. In particular, the map $\tau_\alpha: G \to Z(G)$ defined by
% \[
%   \tau_\alpha(x) = x^{-1}\alpha(x)
% \]
% is a derivation (i.e. 1-cocycle) which is easy to check as $Z(G)$ is abelian and hence these are even
% homomorphisms (see~\cite{NateGaloisCohom} for more details). Since $Z(G)$ is abelian and the action is
% trivial, this derivation is simply a homomorphism from $G$ to $Z(G)$. Since the target is abelian,
% it factors through the commutator subgroup $G'$. Thus, we have an embedding:
% \[
%     \Aut_c(G) \hookrightarrow \Hom(G/G', Z(G))
% \]
% The condition $\Aut(G) = \Aut_c(G)$ implies that \emph{every} automorphism of $G$ acts trivially on the
% quotient $Q = G/Z(G)$. This gives an interesting short exact sequence:
% \[
%   0 \to \Hom(G/G', Z(G)) \to \Aut_c(G) \to \Inn(G) \cap Z(\Aut(G))
% \]
%
%
% Now, let $Q = G/Z(G)$. This notation is used as when we were working with $\Aut(G) = \Q$ we had that $Q$ was
% locally cyclic. In our case, this need not be the case, but it is still divisible and hence $Q$ is suggestive
% of the useful qualities it posses. As we are assuming $G$ is non-abelian, $G/Z(G)$ is non-trivial and:
% \[
%   \omega: Q\times Q \to Z(G) \qquad (a,b) \mapsto [a,b]
% \]
% is non-trivial (if it were trivial, then for any choice of representative $xZ(G),yZ(G)$ of cosets we would get that
% $[x,y] = e$ which would give that $G$ is abelian).


With this, we define a specific map $\Psi: G \to G$ and show it is an automorphism that violates the
constraint established.
Let $Q = G/Z(G)$ and $Z = Z(G)$. Since $G$ is a central extension of $Z$ by $Q$, elements of $G$ can be identified with pairs $(a, z)$ where $a \in Q$ and $z \in Z$, with the multiplication rule:
\[
(a, z_1) \cdot (b, z_2) = (a+b, z_1 + z_2 + f(a, b))
\]
According to the theory of nilpotent groups (see Warfield, \textit{Nilpotent Groups}, Theorem 5.4), for
torsion-free groups, the factor set $f$ corresponds to an alternating bilinear form (the commutator map).
Bilinearity implies:
\[
f(-a, -b) = (-1)(-1)f(a, b) = f(a, b)
\]
We define the map $\Psi: G \to G$ by:
\[
\Psi((a, z)) = (-a, z)
\]

But now, we have established two contradictory facts about $\Psi$:
\begin{enumerate}
    \item  Since $\Psi \in \operatorname{Aut}(G)$ and $\operatorname{Aut}(G)$ is abelian, $\Psi$ must be a central automorphism. It must induce the identity map on $Q$.
    \[ \Psi(a, z) \equiv (a, z) \pmod Z \implies -a = a \]
    \item  By definition, $\Psi$ induces the inversion map on $Q$.
    \[ \Psi(a, z) = (-a, z) \]
\end{enumerate}
These conditions require $-a = a$, or $2a = 0$, for all $a \in Q$.
Since $Q$ is a subgroup of $\mathbb{Q}^n$, it is \emph{torsion-free}. The only solution to $2a=0$ is $a=0$.
This implies $Q$ is the trivial group, so $G = Z(G)$, and $G$ is abelian. The abelian case was already
established, hence no such group can exist.


\end{document}
